
\subsection*{7.15. On fait bouger \(S\).}
Soit \(S'\) un deuxième trait et
\[
  S'\longrightarrow S
\]
un morphisme de traits, i.e. un morphisme local dominant (i.e. transformant point fermé en point fermé, point générique en point générique). On désigne par \(s'\) (resp. \(t'\)) le point fermé (resp. le point générique) de \(S'\), et on pose \(X'=X\times_SS'\). Supposons \(X'\) noethérien (p.ex. \(X\) de type fini sur \(S\)). On a un diagramme commutatif
\[
\begin{tikzcd}
X & X_t \arrow[l,hook,"j_X"'] & & \\
& & X' \arrow[ull] & X'_{t'} \arrow[l,hook,"i'_{X'}"] \arrow[ull] \\
X_s \arrow[uu,hook,"i_X"] & & & \\
& & X'_{s'} \arrow[uu,hook,"j_{X'}"'] \arrow[ull] &
\end{tikzcd}
\;,
\]
à flèches obliques des morphismes plats. On peut donc prendre les homomorphismes images inverses pour \(K_\bullet\) par toutes les flèches du diagramme (IV 2.12). Appliquant la définition 7.3, on trouve que le diagramme suivant
\begin{equation}\tag{7.15.1}
\begin{tikzcd}[column sep=large,row sep=large]
K_\bullet(X_t) \arrow[r,"\sigma_X"] \arrow[d] &
K_\bullet(X_s) \arrow[d]\\
K_\bullet(X'_{t'}) \arrow[r,"\sigma_{X'}"] &
K_\bullet(X'_{s'})
\end{tikzcd}
\end{equation}
est commutatif, où les flèches verticales sont les homomorphismes images inverses. On en déduit aussitôt un diagramme commutatif analogue, avec \(\Gr_\bullet\) au lieu de \(K_\bullet\)~; et lorsque \(X_s\) et \(X_t\) sont réguliers, on peut écrire encore \(K^\bullet\) et \(\Gr^\bullet\) au lieu de \(K_\bullet\) et \(\Gr_\bullet\)~.

